\section*{第 \ref{ch:b3:clt} 章 --- 特征
函数与中心极限定理}

\begin{solution}{exo:b3:clt:1}
$\intcc{-1}1$ 上的均匀分布：$\varphi(\xi) =
\frac12\int_{-1}^1\eu^{\iu\xi x}\dd x =
\frac{\sin\xi}{\xi}$（$1$ 为其在 $\xi = 0$ 处的值）。
指数分布 $\mathcal E(\lambda)$：$\varphi(\xi) =
\lambda\int_0^\infty\eu^{(\iu\xi - \lambda)x}\dd x =
\frac{\lambda}{\lambda - \iu\xi}$（原函数
在 $+\infty$ 处趋于零，因为 $\operatorname{Re}(\iu\xi -
\lambda) < 0$）。泊松分布 $\mathcal P(\lambda)$：由
离散分布的变换定理，
\[
\varphi(\xi) = \sum_{k\geq0}\eu^{\iu\xi
k}\,\eu^{-\lambda}\frac{\lambda^k}{k!}
= \eu^{-\lambda}\exp\bigl(\lambda\eu^{\iu\xi}\bigr)
= \exp\bigl(\lambda(\eu^{\iu\xi} - 1)\bigr).
\]
二项分布 $\mathcal B(n, p)$：$n$ 个独立
伯努利变量之和，每个的特征函数为 $1 - p + p\eu^{\iu\xi}$，
故 $\varphi(\xi) = \bigl(1 - p + p\eu^{\iu\xi}\bigr)^n$
（由\cref{prop:b3:clt:cfbasics}(b)）。泊松分布的可加性：若 $X
\sim \mathcal P(\lambda)$，$Y \sim \mathcal P(\mu)$
相互独立，
\[
\varphi_{X+Y}(\xi) = \eu^{\lambda(\eu^{\iu\xi}-1)}
\eu^{\mu(\eu^{\iu\xi}-1)} =
\eu^{(\lambda+\mu)(\eu^{\iu\xi}-1)},
\]
即 $\mathcal P(\lambda + \mu)$ 的特征函数；由单射性
（\cref{thm:b3:clt:injectivity}）可确定其分布。
\end{solution}

\begin{solution}{exo:b3:clt:2}
(a) $\overline{\varphi_X(\xi)} = \E\eu^{-\iu\xi X} =
\varphi_{-X}(\xi)$。于是 $\varphi_X$ 为实数，当且仅当
$\varphi_X = \varphi_{-X}$，当且仅当（由单射性，
\cref{thm:b3:clt:injectivity}）$X$ 与 $-X$ 具有相同
分布。
(b) 记 $\varphi_X(\xi_0) = \eu^{\iu\theta}$。于是
\[
\E\bigl[1 - \cos(\xi_0X - \theta)\bigr]
= 1 - \operatorname{Re}\bigl(\eu^{-\iu\theta}
\varphi_X(\xi_0)\bigr) = 1 - 1 = 0 .
\]
被积函数非负，故 $\cos(\xi_0X - \theta) = 1$
几乎必然成立（非负随机变量的期望为零时，
它几乎必然为零），即 $\xi_0X - \theta \in 2\pi\Z$ 几乎必然成立：$X$
取值于等差数列
$\frac{\theta}{\xi_0} + \frac{2\pi}{\xi_0}\Z$ 中，几乎必然。
若 $X$ 有密度，则此可数集为勒贝格零测集，
从而其概率为 $0$ --- 矛盾；因此
$\abs{\varphi_X(\xi)} < 1$，对每个 $\xi \neq 0$ 成立。
\end{solution}

\begin{solution}{exo:b3:clt:3}
由独立性及\cref{prop:b3:clt:cfbasics}：
\[
\varphi_{X+Y}(\xi) = \eu^{\iu m_1\xi - \sigma_1^2\xi^2/2}\,
\eu^{\iu m_2\xi - \sigma_2^2\xi^2/2}
= \eu^{\iu(m_1+m_2)\xi - (\sigma_1^2+\sigma_2^2)\xi^2/2},
\]
即 $\mathcal N(m_1 + m_2, \sigma_1^2 + \sigma_2^2)$ 的特征函数，
由单射性即得。仿射变换下的稳定性即为
仿射法则（$aX + b \sim \mathcal N(am_1 + b,
a^2\sigma_1^2)$，允许退化情形 $a = 0$），并
独立和下的稳定性可对
上述计算用归纳法得出。对\emph{相依}的高斯变量，其和
未必为高斯；见\cref{exo:b3:clt:9}，$X$ 与 $Y =
\varepsilon X$ 各自为标准高斯变量，但 $X + Y$
以概率 $\frac12$ 取零，却并非几乎必然
为零，因此它不是高斯变量。
\end{solution}

\begin{solution}{exo:b3:clt:4}
(a) \emph{直接方向。}设 $t$ 为
$F_X$ 的连续点，$\delta > 0$。取连续斜坡函数
$f^-$（$= 1$ 于 $\intoc{-\infty}{t-\delta}$ 上，$0$ 自 $t$ 起，
中间线性）与 $f^+$（$= 1$ 于 $\intoc{-\infty}t$ 上，
$0$ 自 $t + \delta$ 起，中间线性）；于是 $f^- \leq
\mathbf 1_{\intoc{-\infty}t} \leq f^+$，故
\[
\E f^-(X_n) \leq F_{X_n}(t) \leq \E f^+(X_n),
\]
两端项收敛于 $\E f^\pm(X)$，它们本身
又夹在 $F_X(t - \delta)$ 与 $F_X(t + \delta)$ 之间。
令 $n \to \infty$，再令 $\delta \to 0$，并利用
$F_X$ 在 $t$ 处的连续性，得 $F_{X_n}(t) \to F_X(t)$。

\emph{逆方向。}设 $f$ 为有界连续函数，$M =
\sup\abs f$，$\varepsilon > 0$。$F_X$ 的连续点
稠密（$F_X$ 至多有可数个跳跃点），因此
可取连续点 $a < b$，使 $F_X(a) < \varepsilon$
且 $1 - F_X(b) < \varepsilon$。在紧区间 $\intcc ab$ 上，
函数 $f$ 一致连续，可取
$a = t_0 < t_1 < \dots < t_m = b$ 为 $F_X$ 的连续点，使
$f$ 在每个区间上的振幅至多为 $\varepsilon$，
其中区间为 $\intoc{t_{j-1}}{t_j}$，并令 $g = \sum_j
f(t_j)\,\mathbf 1_{\intoc{t_{j-1}}{t_j}}$。于是 $\abs{f - g}
\leq \varepsilon$（在 $\intoc ab$ 上），$\abs g \leq M$，且对
$T = X_n$ 或 $X$：
\[
\bigl|\E f(T) - \E g(T)\bigr| \leq \varepsilon +
2M\bigl(F_T(a) + 1 - F_T(b)\bigr).
\]
此外 $\E g(X_n) = \sum_j f(t_j)\bigl(F_{X_n}(t_j) -
F_{X_n}(t_{j-1})\bigr) \to \E g(X)$（有限个收敛项之和，
所有 $t_j$ 均为连续点），并且
$F_{X_n}(a) \to F_X(a) < \varepsilon$，$1 - F_{X_n}(b) \to 1
- F_X(b) < \varepsilon$。综合可得：
$\limsup_n\abs{\E f(X_n) - \E f(X)} \leq 2\varepsilon +
8M\varepsilon$；令 $\varepsilon \to 0$。

(b) 常数 $c$ 的分布函数为
$\mathbf 1_{\intco c\infty}$，除 $c$ 外连续。对于
$\varepsilon > 0$，点 $c - \varepsilon$ 与 $c +
\frac\varepsilon2$ 均为连续点，故
\[
\P(\abs{X_n - c} > \varepsilon) \leq F_{X_n}(c -
\varepsilon) + 1 - F_{X_n}\Bigl(c + \frac\varepsilon2\Bigr)
\longrightarrow 0 + 1 - 1 = 0 .
\]
\end{solution}

\begin{solution}{exo:b3:clt:5}
设 $z_n = p_n(\eu^{\iu\xi} - 1)$，则 $\varphi_{X_n}(\xi) =
(1 + z_n)^n$（见\cref{exo:b3:clt:1}），且 $\abs{z_n} \leq 2p_n
\to 0$（注意 $p_n = \frac{np_n}n \to 0$）。$1 + z_n$ 与
$\eu^{z_n}$ 的模均不超过 $1$：$\abs{1 + z_n} =
\abs{(1 - p_n) + p_n\eu^{\iu\xi}} \leq 1$，由三角
不等式可得，且 $\abs{\eu^{z_n}} = \eu^{p_n(\cos\xi - 1)}
\leq 1$。伸缩不等式 $\abs{a^n - b^n} \leq
n\abs{a - b}$（\cref{thm:b3:clt:clt} 的证明）以及幂
级数估计 $\abs{\eu^z - 1 - z} \leq
\abs z^2\eu^{\abs z}$ 给出
\[
\bigl|(1 + z_n)^n - \eu^{nz_n}\bigr| \leq n\bigl|1 + z_n -
\eu^{z_n}\bigr| \leq n\,\abs{z_n}^2\,\eu^{\abs{z_n}} \leq
4\eu^2\,np_n^2 = 4\eu^2\,(np_n)\,p_n \longrightarrow 0 .
\]
由 $nz_n = np_n(\eu^{\iu\xi} - 1) \to
\lambda(\eu^{\iu\xi} - 1)$，可得 $\varphi_{X_n}(\xi)
\to \exp\bigl(\lambda(\eu^{\iu\xi} - 1)\bigr)$，对每个
$\xi$ 成立，即 $\mathcal P(\lambda)$ 的特征函数，由列维
连续性定理（\cref{thm:b3:clt:levy}）可得 $X_n \Rightarrow \mathcal
P(\lambda)$。数值上：$\P\bigl(\mathcal B(100, 0.02) =
0\bigr) = 0.98^{100} = \eu^{100\ln 0.98} \approx
\eu^{-2.020} \approx 0.1326$，而 $\P\bigl(\mathcal P(2) =
0\bigr) = \eu^{-2} \approx 0.1353$：两者仅相差约百分之二，
已在这样粗略的 $n$ 下如此。
\end{solution}

\begin{solution}{exo:b3:clt:6}
(a) 掷一次骰子的均值为 $\frac72$，方差
为 $\frac{35}{12}$，故 $S$ 的均值为 $3500$，方差
为 $\frac{35000}{12} \approx 2916.7$，标准差
为 $\approx 54.0$。由中心极限定理，
\[
\P(S > 3600) = \P\Bigl(\frac{S - 3500}{54.0} > 1.85\Bigr)
\approx 1 - \Phi(1.85) \approx 0.032 :
\]
约 $3\%$ 的概率。
(b) $S \sim \mathcal B(100, \frac12)$：均值为 $50$，标准
差为 $5$。利用连续性修正，
\[
\P(45 \leq S \leq 55) \approx
\Phi\Bigl(\frac{55.5 - 50}{5}\Bigr) -
\Phi\Bigl(\frac{44.5 - 50}{5}\Bigr) = 2\Phi(1.1) - 1
\approx 0.729,
\]
与精确值 $0.7287$ 相比；若不作修正，
$2\Phi(1) - 1 \approx 0.683$，偏差将近五个百分点。修正
之所以重要，是因为 $S$ 为格点变量：每个
原子 $\P(S = k)$ 可由
$\intcc{k - \frac12}{k + \frac12}$ 上的高斯质量很好地逼近，而将
区间在整数 $45$ 与 $55$ 处截断时，
两端各丢掉半个原子。
\end{solution}

\begin{solution}{exo:b3:clt:7}
变量 $g(U_k)$ 独立同分布（它们是
独立同分布变量的可测像），平方可积，均值为 $I$
（由变换定理，\cref{exo:b3:product:9}）及方差
为 $\sigma^2$。若 $\sigma > 0$，则 \cref{thm:b3:clt:clt} 作用于
它们，恰好得到所述收敛：
\[
\sqrt n\,\Bigl(\frac1n\sum_{k\leq n}g(U_k) - I\Bigr) =
\frac{\sum_{k\leq n}\bigl(g(U_k) - I\bigr)}{\sqrt n}
\Longrightarrow \mathcal N(0, \sigma^2)
\]
（若 $\sigma = 0$，则 $g$ 几乎处处为常数，于是左端
恒为零）。由此可得
$\P\bigl(\abs{\frac1n\sum g(U_k) - I} \leq
1.96\,\sigma/\sqrt n\bigr) \to 0.95$：误差棒 $\pm
1.96\,\sigma/\sqrt n$ 对维数 $d$ 的感知仅
经由常数 $\sigma$，而绝不经由 $n$ 的速率。
中点法则取 $n$ 个节点、维数为 $d$ 时，网格宽为
$n^{-1/d}$，误差阶为 $n^{-2/d}$，对 $\mathcal C^2$
被积函数。蒙特卡罗方法的 $n^{-1/2}$ 衰减
比 $n^{-2/d}$ 更快，恰在 $\frac12 > \frac2d$，即 $d > 4$：
自 $5$ 维起，随机抽样渐近地优于
网格 --- 维数灾难不会困扰
概率方法，这正是蒙特卡罗方法主导
高维积分的原因。
\end{solution}

\begin{solution}{exo:b3:clt:8}
\emph{和。}对固定的 $\xi$：
\[
\bigl|\E\eu^{\iu\xi(X_n+Y_n)} -
\eu^{\iu\xi c}\,\E\eu^{\iu\xi X_n}\bigr|
= \bigl|\E\bigl[\eu^{\iu\xi X_n}\bigl(\eu^{\iu\xi Y_n} -
\eu^{\iu\xi c}\bigr)\bigr]\bigr|
\leq \E\bigl|\eu^{\iu\xi(Y_n - c)} - 1\bigr| .
\]
按事件 $\{\abs{Y_n - c} \leq \delta\}$ 分情形讨论：在该事件上，
$\abs{\eu^{\iu\xi(Y_n-c)} - 1} \leq \abs\xi\,\delta$（弦
比弧短）；补集上的贡献
至多为 $2\,\P(\abs{Y_n - c} > \delta) \to 0$。由此
$\limsup$ 满足 $\leq \abs\xi\,\delta$，对每个 $\delta > 0$ 成立：
差值趋于 $0$。由于 $\E\eu^{\iu\xi X_n} \to
\varphi_X(\xi)$，故 $\varphi_{X_n+Y_n}(\xi) \to
\eu^{\iu\xi c}\varphi_X(\xi) = \varphi_{X+c}(\xi)$，而由
列维连续性定理（\cref{thm:b3:clt:levy}）可得 $X_n + Y_n
\Rightarrow X + c$。

\emph{乘积。}首先，$cX_n \Rightarrow cX$：
$\varphi_{cX_n}(\xi) = \varphi_{X_n}(c\xi) \to
\varphi_X(c\xi) = \varphi_{cX}(\xi)$。其次，$(Y_n - c)X_n
\to 0$ 依概率成立：$X_n$ 的分布是紧的（它们的
特征函数收敛于某个特征函数；参见
\cref{thm:b3:clt:levy}中的紧性步骤），故对给定的 $\varepsilon > 0$，取 $M$
使得 $\P(\abs{X_n} > M) \leq \varepsilon$ 对所有 $n$ 成立；则
\[
\P\bigl(\abs{(Y_n - c)X_n} > \varepsilon\bigr) \leq
\P(\abs{X_n} > M) + \P\Bigl(\abs{Y_n - c} >
\frac{\varepsilon}{M}\Bigr) \leq \varepsilon + o(1) .
\]
写成 $Y_nX_n = cX_n + (Y_n - c)X_n$，并应用和的部分
（其证明只用到 $Y_n' := (Y_n - c)X_n \to 0$ 依概率
收敛，且常数为 $0$）：$Y_nX_n \Rightarrow cX$。

\emph{应用。}由强大数定律
（\cref{thm:b3:probability:slln}），$\hat p_n \to p$ 几乎必然，故
由连续性 $\hat\sigma_n = \sqrt{\hat p_n(1 - \hat p_n)}
\to \sigma = \sqrt{p(1 - p)} > 0$ 几乎必然，从而
$\frac{\sigma}{\hat\sigma_n} \to 1$ 依概率成立。斯卢茨基
乘积规则将 $\frac{S_n - np}{\sigma\sqrt n}
\Rightarrow \mathcal N(0,1)$ 升级为 $\frac{S_n -
np}{\hat\sigma_n\sqrt n} = \frac{\sigma}{\hat\sigma_n}\cdot
\frac{S_n - np}{\sigma\sqrt n} \Rightarrow \mathcal N(0,1)$：
即\emph{可用的}置信区间 $\hat p_n \pm
1.96\,\hat\sigma_n/\sqrt n$，仅由数据构造，
其渐近 $95\%$ 水平保持不变。
\end{solution}

\begin{solution}{exo:b3:clt:9}
(a) 按
$\varepsilon$ 的两个取值拆分期望（由独立性）：对博雷尔集 $B$，$\P(Y \in B) =
\frac12\P(X \in B) + \frac12\P(-X \in B) = \P(X \in B)$，
由于 $-X \sim X$（$\mathcal N(0,1)$ 对称）：$Y \sim
\mathcal N(0,1)$。又 $\operatorname{Cov}(X, Y) =
\E[\varepsilon X^2] = \E[\varepsilon]\,\E[X^2] = 0 \cdot 1 =
0$。
(b) $\abs Y = \abs X$，故 $\P(\abs X \leq 1,\ \abs Y \geq 2)
= 0$，而 $\P(\abs X \leq 1)\,\P(\abs Y \geq 2) > 0$：不
独立。若 $(X, Y)$ 为高斯向量，则 $X + Y =
(1 + \varepsilon)X$ 将是一个实高斯变量
（见\cref{def:b3:clt:gaussianvector}，取 $t = (1,1)$）；但
$\P(X + Y = 0) = \P(\varepsilon = -1) = \frac12$，而
高斯变量仅当几乎必然为常数时才有原子
--- 并且 $X + Y$ 几乎必然等于 $2X \neq 0$，
在 $\{\varepsilon = 1\}$ 上成立。矛盾：$(X, Y)$ 不是
高斯向量。
(c) 每个边缘分布都是高斯的，且协方差为零，
但独立性不成立 --- 因为\emph{二元组}不是
联合高斯。\cref{thm:b3:clt:gaussianvector}(2) 不能
减弱为“边缘均为高斯”。
\end{solution}

\begin{solution}{exo:b3:clt:10}
(a) 由\cref{exo:b3:fouriertransform:1}算得
$\widehat{\eu^{-\abs\cdot}}(\xi) = \frac{2}{1 + \xi^2}$；
两边均可积，由傅里叶反演
（\cref{thm:b3:fouriertransform:inversion}）将此式
反过来：
\[
\int_\R\eu^{\iu\xi x}\,\frac{\dd x}{\pi(1 + x^2)} =
\eu^{-\abs\xi},
\]
这恰是 $\varphi_X(\xi)$，即柯西随机变量 $X$ 的特征函数。
(b) 由独立性，$\varphi_{S_n}(\xi) =
\bigl(\eu^{-\abs\xi}\bigr)^n = \eu^{-n\abs\xi}$，故
$\varphi_{S_n/n}(\xi) = \varphi_{S_n}(\xi/n) =
\eu^{-\abs\xi}$：可见经验均值 $\frac{S_n}n$
对每个 $n$ 仍服从标准柯西分布（由单射性）。平均值
永不集中：在 $10^6$ 时，其波动
与单次观测的波动相同。
(c) $\E\abs{X_1} = \frac2\pi\int_0^\infty\frac{x\,\dd x}{1 +
x^2} = +\infty$：柯西分布不可积，故
强大数定律
（\cref{thm:b3:probability:slln}）不适用，而中心极限定理
（需要有限方差）更不适用。此处它们的
结论真正失效，而不仅仅是证明失效。
一致性检验：$\varphi(\xi) = \eu^{-\abs\xi}$ 在
$0$ 处不可导，正如\cref{prop:b3:clt:cfbasics}(c)
逆否读出时对不可积的
随机变量所预言的那样。
\end{solution}

\begin{solution}{exo:b3:clt:11}
(a) $\varphi_{aX_1 + bX_2}(\xi) =
\eu^{-a\abs\xi}\eu^{-b\abs\xi} = \eu^{-(a+b)\abs\xi} =
\varphi_{(a+b)X_1}(\xi)$（由独立性及
\cref{exo:b3:clt:10}），由单射性知二者分布相同。

(b) $\varphi_{aX_1+bX_2}(\xi) = \eu^{-a^2\xi^2/2}
\eu^{-b^2\xi^2/2} = \eu^{-(a^2+b^2)\xi^2/2}$：即
$\sqrt{a^2+b^2}\,X_1$ 的分布。

(c) 若 $\varphi_X(\xi) = \eu^{-c\abs\xi^\alpha}$，则
$S_n = X_1 + \dots + X_n$ 的特征函数为 $\varphi_{S_n} =
\eu^{-cn\abs\xi^\alpha}$，而 $S_n/n^{1/\alpha}$ 的特征函数为
$\varphi(\xi) = \eu^{-c\abs\xi^\alpha}$，仍
是精确自再生，在 $n^{1/\alpha}$ 的尺度下 ---
对高斯情形为 $\sqrt n$（$\alpha = 2$），而
柯西情形则为 $n$ 本身（$\alpha = 1$，见\cref{exo:b3:clt:10}(b)）。独立同分布
的变量之和（经仿射
标准化）只能收敛到一个在此类卷积下稳定的
分布；中心极限定理指出有限方差迫使其落入
高斯吸引域，而柯西吸引域留给
尾部极重的分布，以至 $\E X^2 = \infty$ 甚至 $\E\abs
X = \infty$ --- 例如柯西变量本身的和。
普适性有若干“孤岛”，按尾部
指数 $\alpha \in \intoc02$ 区分。
\end{solution}

\begin{solution}{exo:b3:clt:12}
(a) 指示变量 $\mathbf 1_{X_k \leq t}$ 独立同分布，
服从参数 $p = F(t)$ 的伯努利分布：其和 $nF_n(t)$
服从二项分布 $\mathcal B(n, p)$；由强大数定律，$F_n(t)
\to p$ 几乎必然成立，对同一指示变量应用中心极限定理（
\cref{thm:b3:clt:clt}；方差 $p(1-p)$）得所述
高斯极限。

(b) $p(1 - p)$ 在 $p = \frac12$ 处取最大值，亦即
$F(t) = \frac12$ 时，即\emph{中位数}处。估计尾概率
渐近地较容易（方差 $\to 0$，当
$p \to 0, 1$ 时）；中位数附近承载着最大的
统计噪声 --- 经验曲线在
中部波动最剧烈。

(c) 设 $F$ 连续，则变量 $U_k = F(X_k)$
独立同分布，且在 $\intoo01$ 上均匀分布（见\cref{exo:b3:probability:1}），
由 $F$ 的单调性，记 $G_n$ 为
$U_k$ 的经验分布函数，有
\[
\sup_{t\in\R}\,\abs{F_n(t) - F(t)}
= \sup_{u \in \operatorname{im}F}\,\abs{G_n(u) - u}
= \sup_{u\in\intcc01}\abs{G_n(u) - u} :
\]
第一个等号成立，是因为 $\{X_k \leq t\} = \{U_k \leq
F(t)\}$ 除零概率事件外（由单调性；严格不等号
只可能在 $F$ 的平坦段失效，而那里两边的值
都不变），第二个等号则因为连续的 $F$，
从 $0$ 增至 $1$，取遍
$\intoo01$ 中的每个值（由介值定理），端点
无贡献（$G_n(0) - 0 = 0$ 与 $G_n(1) - 1 = 0$）。右端
只涉及均匀变量：分布对一切
$F$ 都相同 --- 故只需一张临界值表（即柯尔莫哥洛夫
分布的临界值表）即可检验\emph{任一}连续
模型与数据是否相符。
\end{solution}

\begin{solution}{pb:b3:clt:1}
\textbf{1.} 族 $(X_1, \dots, X_n, N_1, \dots, N_n)$
是独立的：两个块彼此独立
（按构造），且每个块都独立同分布。$W_i$ 是
仅由变量 $(X_j)_{j<i}$ 与
$(N_j)_{j>i}$ 决定的可测函数，这些变量均不同于 $X_i$ 与 $N_i$。
由分组原理（
\cref{thm:b3:probability:independence}），$W_i$
与 $(X_i, N_i)$ 独立。两式
$H_i = W_i + \frac{X_i}{\sqrt n}$ 与 $H_{i-1} = W_i +
\frac{N_i}{\sqrt n}$ 由定义直接可得：
从 $H_i$ 到 $H_{i-1}$，只需把单个加项
$X_i$ 换为 $N_i$。

\textbf{2.} 泰勒–拉格朗日公式（$3$ 阶）：存在 $c$
介于 $w$ 与 $w + h$ 之间，使得 $f(w + h) = f(w) + f'(w)h +
\frac12f''(w)h^2 + \frac16f'''(c)h^3$，且 $\abs{f'''(c)}
\leq M_3$ 即得所求界。

\textbf{3.} 在公共
基点 $w = W_i$ 处将两式相减：
\[
f(H_i) - f(H_{i-1}) = f'(W_i)\,\frac{X_i - N_i}{\sqrt n} +
\frac{f''(W_i)}{2}\,\frac{X_i^2 - N_i^2}{n} + R_i,
\qquad
\abs{R_i} \leq \frac{M_3}{6}\cdot
\frac{\abs{X_i}^3 + \abs{N_i}^3}{n^{3/2}} .
\]
两边取期望。由第 1 问，$f'(W_i)$ 及 $f''(W_i)$
均与 $(X_i, N_i)$ 独立，故混合期望
可以分解：
\begin{align*}
\E\Bigl[f'(W_i)\,\frac{X_i - N_i}{\sqrt n}\Bigr] &=
\E\bigl[f'(W_i)\bigr]\,\frac{\E X_i - \E N_i}{\sqrt n} = 0,
\\
\E\Bigl[f''(W_i)\,\frac{X_i^2 - N_i^2}{n}\Bigr] &=
\E\bigl[f''(W_i)\bigr]\,\frac{1 - 1}{n} = 0 :
\end{align*}
$X_i$ 与 $N_i$ 的前两阶矩\emph{相等}，
故只有余项保留：
\[
\bigl|\E f(H_i) - \E f(H_{i-1})\bigr| \leq \E\abs{R_i} \leq
\frac{M_3}{6}\cdot\frac{\beta + \gamma}{n^{3/2}} .
\]
高斯分布的三阶矩：$\gamma = \E\abs{N_1}^3 =
2\int_0^\infty x^3\,\frac{\eu^{-x^2/2}}{\sqrt{2\pi}}\,\dd x
= \frac{2}{\sqrt{2\pi}}\int_0^\infty 2u\,\eu^{-u}\dd u =
\frac{4}{\sqrt{2\pi}} = \frac{2\sqrt2}{\sqrt\pi}$
（换元 $u = x^2/2$，再由 $\Gamma(2) = 1$）。

\textbf{4.} 逐项相消：$\E f(T_n) - \E f(G_n) =
\sum_{i=1}^n\bigl(\E f(H_i) - \E f(H_{i-1})\bigr)$，并对
这 $n$ 项中的每一项应用第 3 问：
\[
\bigl|\E f(T_n) - \E f(G_n)\bigr| \leq
n \cdot \frac{M_3(\beta + \gamma)}{6\,n^{3/2}} =
\frac{M_3\,(\beta + \gamma)}{6\,\sqrt n} .
\]

\textbf{5.} $G_n$ 服从 $\mathcal N(0,1)$，\emph{精确}对每个
$n$ 成立（独立标准高斯变量的规范化和，
见\cref{exo:b3:clt:3}），故 $\E f(G_n) = \E f(N)$，而
第 4 问化为 $\abs{\E f(T_n) - \E f(N)} \leq
\frac{M_3(\beta+\gamma)}{6\sqrt n} \to 0$，对 $f \in
\mathcal C^3_b$。\emph{推广。}设 $f$ 为有界
连续函数，$M = \sup\abs f$，$\varepsilon > 0$。取 $A
\geq 1$ 使 $\frac1{A^2} \leq \varepsilon$。由切比雪夫不等式及
$\V(T_n) = 1$ 得 $\P(\abs{T_n} > A) \leq \varepsilon$
对所有 $n$ 成立，同理 $\P(\abs N > A) \leq \varepsilon$。
取 $\chi$ 为 $\mathcal C^\infty$ 函数，满足
$\mathbf 1_{\intcc{-A}A} \leq \chi \leq
\mathbf 1_{\intcc{-A-1}{A+1}}$（光滑平台，由
$\mathbf 1_{\intcc{-A-\frac12}{A+\frac12}}$ 磨光而得，
见\cref{thm:b3:lp:regularization}）；$g = f\chi$ 连续，
具紧支集，从而一致连续，故其
磨光 $g_\eta = g * \rho_\eta$ 为 $\mathcal
C^\infty$ 类，各阶导数有界，且
$\norm{g - g_\eta}_\infty \leq \varepsilon$，当 $\eta$
充分小时成立。对 $T = T_n$ 或 $N$，因 $f = g$ 在
$\intcc{-A}A$ 上，且 $\abs{f - g} \leq 2M$ 处处成立：
\[
\bigl|\E f(T) - \E g_\eta(T)\bigr| \leq
\E\abs{(f - g)(T)} + \norm{g - g_\eta}_\infty
\leq 2M\,\P(\abs T > A) + \varepsilon \leq
(2M + 1)\,\varepsilon .
\]
结合 $\E g_\eta(T_n) \to \E g_\eta(N)$（第
4 问适用，因 $g_\eta \in \mathcal C^3_b$）：
\[
\limsup_n\;\bigl|\E f(T_n) - \E f(N)\bigr| \leq
2(2M + 1)\,\varepsilon ,
\]
由 $\varepsilon$ 的任意性，$\E f(T_n) \to \E f(N)$ 对
每个有界连续 $f$ 成立，即 $T_n \Rightarrow
\mathcal N(0,1)$。

\textbf{6.} 同分布只在一句话
中出现：“$X_i$ 与 $N_i$ 具有相同的前两阶
矩”。故设 $X_1, \dots, X_n$ 相互独立，
中心化、方差为 $\sigma_i^2$ 且三阶矩有限，
$s_n^2 = \sum_i\sigma_i^2 > 0$，并取 $N_i \sim
\mathcal N(0, \sigma_i^2)$，且与一切量独立。
以 $s_n$ 为规范化因子定义混合和：$H_i =
\frac1{s_n}(\sum_{j\leq i}X_j + \sum_{j>i}N_j)$。在
第 $i$ 次交换中，$\E X_i = \E N_i = 0$ 且 $\E X_i^2 = \E N_i^2
= \sigma_i^2$ 再次消去 $f'$ 与 $f''$ 项，并
余项给出（利用 $\E\abs{N_i}^3 = \sigma_i^3\gamma$，由
尺度变换）：
\[
\bigl|\E f(H_i) - \E f(H_{i-1})\bigr| \leq
\frac{M_3}{6\,s_n^3}\bigl(\E\abs{X_i}^3 +
\sigma_i^3\gamma\bigr) .
\]
裂项求和：
\[
\Bigl|\E f\Bigl(\frac{X_1 + \dots + X_n}{s_n}\Bigr) -
\E f(N)\Bigr| \leq \frac{M_3}{6\,s_n^3}\sum_{i=1}^n
\Bigl(\E\abs{X_i}^3 + \V(X_i)^{3/2}\,\gamma\Bigr) .
\]
由于 $\sigma_i^3 = (\E X_i^2)^{3/2} \leq \E\abs{X_i}^3$
（幂平均不等式，即对 $t \mapsto
t^{3/2}$ 作用于 $X_i^2$ 的琴生不等式），右端至多为 
$\frac{M_3(1 + \gamma)}{6}\cdot
\frac{\sum_i\E\abs{X_i}^3}{s_n^3}$：在\emph{李雅普诺夫
条件}下 $\frac1{s_n^3}\sum_i\E\abs{X_i}^3 \to 0$，
标准化和依分布收敛于 $\mathcal N(0,1)$ ---
即不要求同分布的中心极限定理。

\textbf{7.} 设 $\rho \in \mathcal C^\infty_c(\intoo01)$
满足 $\int\rho = 1$，并令 $\psi(x) = \int_x^1\rho(s)\dd
s$：$\psi$ 为 $\mathcal C^\infty$ 函数且单调不增，$\psi = 1$
在 $\R_-$ 上，$\psi = 0$ 在 $\intco1\infty$ 上；令 $K =
\norm{\psi'''}_\infty$。对 $t \in \R$ 及 $\delta > 0$，
定义 $\psi_\delta(x) = \psi\bigl(\frac{x - t}\delta\bigr)$
及 $\tilde\psi_\delta(x) = \psi\bigl(\frac{x - t}\delta +
1\bigr)$：它们都属于 $\mathcal C^3_b$，三阶导数
以 $K/\delta^3$ 为界，并且
\[
\mathbf 1_{\intoc{-\infty}{t-\delta}} \leq
\tilde\psi_\delta \leq \mathbf 1_{\intoc{-\infty}t} \leq
\psi_\delta \leq \mathbf 1_{\intoc{-\infty}{t+\delta}} .
\]
上界：由第 4 题对 $\psi_\delta$ 应用（
取 $M_3 = K/\delta^3$），得
\[
\P(T_n \leq t) \leq \E\psi_\delta(T_n) \leq
\E\psi_\delta(N) + \frac{K(\beta +
\gamma)}{6\,\delta^3\sqrt n}
\leq \Phi(t + \delta) + \frac{K(\beta +
\gamma)}{6\,\delta^3\sqrt n}
\leq \Phi(t) + \frac{\delta}{\sqrt{2\pi}} +
\frac{K(\beta + \gamma)}{6\,\delta^3\sqrt n},
\]
因为 $\Phi$ 是利普希茨的，常数为
$\frac1{\sqrt{2\pi}}$（其密度以
$\frac1{\sqrt{2\pi}}$ 为界）。对称的下界估计经
$\tilde\psi_\delta$ 可得如下两项估计
\[
\sup_{t\in\R}\,\bigl|\P(T_n \leq t) - \Phi(t)\bigr| \leq
\frac{K(\beta + \gamma)}{6}\cdot\frac{1}{\delta^3\sqrt n} +
\frac{\delta}{\sqrt{2\pi}}
\qquad(\delta > 0\ \text{任意}).
\]
两项在 $\delta^{-3}n^{-1/2} \asymp
\delta$（即 $\delta = n^{-1/8}$）时平衡，此时二者均为
$O(n^{-1/8})$，对每个 $n$ 都成立的显式一致速率。
（最优的贝里–埃森速率 $C\beta/\sqrt n$
需要傅里叶光滑化不等式；交换法则以
精度换取完全的初等性。）

\textbf{8.} 取 $X_i = 2B_i - 1$（等概率正负号）：中心化、
方差为 $1$，且 $\abs{X_i} = 1$，故 $\beta = 1$。第 7 题
随后显式且一致地界定 $\sup_t\abs{\P(\frac{S_n}{\sqrt n} \leq t) -
\Phi(t)}$，对\emph{每个}有限
$n$ 成立 --- 一个整体性、非渐近的
关于分布函数的陈述。局部估计
\cref{pb:b3:product:1} 第 7 题给出的则是单个原子的
精确渐近，$\P(S_{2n} = 2k) \sim
\frac{\eu^{-k^2/n}}{\sqrt{\pi n}}$：它能分辨
量级为 $n^{-1/2}$ 的概率，远低于第 7 题的
$n^{-1/8}$ 分辨率，但它是逐点的、渐近的
（固定 $n$ 时无显式误差），并且依赖于这一
特定格点律。局部精度与整体
一致性：两种方法互为补充，将
局部估计对 $k \in \intint{a\sqrt n}{b\sqrt n}$ 求和，
即可在区间上恢复棣莫弗–拉普拉斯定理 --- 速率
更优，但仅对这一律成立。

\textbf{9.} 交换论证关于 $X_i$ 的分布，除
$\E X_i = 0$、$\E X_i^2 = 1$ 之外，还用到
$\E\abs{X_i}^3$ 的有限性：若将高斯变量
$N_i$ 换成任意其他独立同分布族，且前两阶矩相同、
三阶矩有限，则同一裂项求和会
给出 $\abs{\E f(\text{和}_X) - \E f(\text{和}_Y)}$ 的
上界 $O(n^{-1/2})$，对每个光滑的 $f$ 成立。大量独立和的光滑统计量
是\emph{普适的}：在
可量化的误差范围内，它们对求和项分布的依赖
仅通过两个数。这就是不变性原理：
先对最易计算的分布证明极限定理（
高斯分布，一切均可精确计算），然后
借助交换法将其推广到所有分布。同一方案 --- 把和式
换成更复杂的泛函 --- 同样推动了维格纳半圆律
对随机矩阵的情形，随机多项式根的普适性，
以及现代概率论的大部分内容；中心
极限定理只是其中最早、最简单的一个实例。

\textbf{10.} $Z$ 仅是 $(N_{i+1}, \dots,
N_n)$ 的博雷尔函数，而 $A = W_i + \theta h - Z$ 与 $h$ 都是
独立族中其余变量的函数。
$(X_1, \dots, X_n, N_1, \dots, N_n)$ 这一族，由
分组原理（\cref{thm:b3:probability:independence}）可知，
$Z$ 与 $(A, h)$ 独立。作为独立的
$N_j/\sqrt n \sim \mathcal N(0, \frac1n)$ 之和，$Z \sim \mathcal
N(0, s^2)$，其中 $s^2 = \frac{n-i}n$（见\cref{exo:b3:clt:3}），
其密度以 $\frac1{s\sqrt{2\pi}}$ 为界。
$((A, h), Z)$ 的分布是两个边缘分布的乘积，故
由托内利定理（转换公式）固定第一块：记 $G(a) =
\E\abs{g(a + Z)} = \int\abs{g(a + z)}\,\varphi_s(z)\,\dd z
\leq \frac{\norm g_{L^1}}{s\sqrt{2\pi}}$ 对每个 $a$，
\[
\E\bigl[\abs h^3\abs{g(A + Z)}\bigr] =
\E\bigl[\abs h^3\,G(A)\bigr] \leq
\frac{\norm g_{L^1}}{s\sqrt{2\pi}}\,\E\abs h^3
= \sqrt{\frac{n}{2\pi(n-i)}}\,\norm g_{L^1}\,\E\abs h^3 .
\]

\textbf{11.} 泰勒公式的积分形式由
$f(w + h) - f(w) =
h\int_0^1f'(w + \theta h)\,\dd\theta$ 关于
$\theta$ 作两次分部积分得到。对第 $i$ 次交换取期望后，
$0, 1, 2$ 阶项恰如第 3 题那样相消，而
两个余项（对 $h = X_i/\sqrt n$
与 $N_i/\sqrt n$）当 $i \leq n - 1$ 时有界，由
第 10 题，取 $g = f'''$ 可得：
\[
\bigl|\E f(H_i) - \E f(H_{i-1})\bigr| \leq
\int_0^1\frac{(1-\theta)^2}2\,\dd\theta\;
\sqrt{\frac{n}{2\pi(n-i)}}\,\norm{f'''}_{L^1}
\frac{\beta + \gamma}{n^{3/2}}
= \frac{\beta + \gamma}{6\,n^{3/2}}
\sqrt{\frac{n}{2\pi(n-i)}}\,\norm{f'''}_{L^1} .
\]
求和时，$\sum_{i=1}^{n-1}\sqrt{\frac n{n-i}} =
\sqrt n\sum_{m=1}^{n-1}m^{-1/2} \leq 2n$，并加上
第 3 题对最后一次交换（$i = n$，此时已无
高斯项）给出的界：
\[
\bigl|\E f(T_n) - \E f(G_n)\bigr| \leq
\frac{\beta + \gamma}{3\sqrt{2\pi}}\cdot
\frac{\norm{f'''}_{L^1}}{\sqrt n} +
\frac{M_3(\beta + \gamma)}{6\,n^{3/2}} .
\]
斜坡函数：$\psi_\delta'''(x) =
\delta^{-3}\psi'''\bigl(\frac{x - t}\delta\bigr)$，故
$M_3 = K\delta^{-3}$，其中 $K = \norm{\psi'''}_\infty$，且
$\norm{\psi_\delta'''}_{L^1} = \delta^{-2}
\norm{\psi'''}_{L^1} = K_1\delta^{-2}$（换元）。
于是由第 7 题的夹逼式给出
\[
\sup_t\,\bigl|\P(T_n \leq t) - \Phi(t)\bigr| \leq
\frac{K_1(\beta + \gamma)}{3\sqrt{2\pi}}\cdot
\frac1{\delta^2\sqrt n} +
\frac{K(\beta + \gamma)}{6}\cdot
\frac1{\delta^3n^{3/2}} + \frac\delta{\sqrt{2\pi}} .
\]
取 $\delta = n^{-1/6}$ 时，第一项与第三项均为
$O(n^{-1/6})$，中间一项为 $O(n^{-1})$：一致
速率 $O(n^{-1/6})$，严格优于第 7 题的
$n^{-1/8}$ --- 混合方案中的高斯部分完成了
额外的光滑化。

\textbf{12.} $\E N_1^3 = 0$（奇被积函数），且
由分部积分得 $\E N_1^4 = 3\,\E N_1^2 = 3$
($\int x^3\cdot x\varphi(x)\dd x =
3\int x^2\varphi$）。设 $f$ 为 $\mathcal C^4$ 类函数，
其导数有界，将每次交换展开到四阶，其中
三阶项含有因子 $\E X_i^3 - \E N_i^3 =
0$（由独立性如第 3 题那样分解），故仅余下
四阶余项
$\int_0^1\frac{(1-\theta)^3}6f^{(4)}(w + \theta
h)h^4\dd\theta$ 得以保留，其中 $\int_0^1
\frac{(1-\theta)^3}6\dd\theta = \frac1{24}$，且 $\E h^4 =
\beta_4n^{-2}$ 或 $3n^{-2}$。由第 10 题（取 $g =
f^{(4)}$）给出 $i \leq n - 1$ 时交换的界，并
仿第 11 题求和：
\[
\bigl|\E f(T_n) - \E f(G_n)\bigr| \leq
\frac{\beta_4 + 3}{12\sqrt{2\pi}}\cdot
\frac{\norm{f^{(4)}}_{L^1}}{n} +
\frac{M_4(\beta_4 + 3)}{24\,n^2} .
\]
设 $\norm{\psi_\delta^{(4)}}_{L^1} = K_2\delta^{-3}$，且
$M_4 = K'\delta^{-4}$，则分布函数的界
变为 $C\delta^{-3}n^{-1} + C'\delta^{-4}n^{-2} +
\frac\delta{\sqrt{2\pi}}$；取 $\delta = n^{-1/4}$ 时，两端项
为 $O(n^{-1/4})$，中间项为 $O(n^{-1})$：
速率 $O(n^{-1/4})$。

\textbf{13.} 在 $k$ 个矩匹配的条件下，每次交换保留的
余项量级为 $\E\abs h^{k+1} \asymp
n^{-(k+1)/2}$；隐含高斯的界以
$\norm{f^{(k+1)}}_{L^1}$ 计入，对交换求和贡献
因子 $2n$，给出 $\asymp\norm{f^{(k+1)}}_{L^1}\,
n^{-(k-1)/2}$（对光滑的 $f$）。斜坡函数的代价为
$\norm{\psi_\delta^{(k+1)}}_{L^1} \asymp \delta^{-k}$，故
分布函数的误差为 $\asymp
\delta^{-k}n^{-(k-1)/2} + \delta$，平衡点为 $\delta =
n^{-(k-1)/(2k+2)}$：速率 $n^{-(k-1)/(2k+2)}$，
即 $n^{-1/6}$（$k = 2$ 时），$n^{-1/4}$（$k = 3$ 时），并趋向
于 $n^{-1/2}$，仅当 $k \to \infty$ 时 --- 但 $k \geq 4$
将迫使 $\E X_1^4 = 3$ 乃至更高矩，亦即一个
分布已在模仿高斯分布。饱和现象是
结构性的：交换使 $n$ 个交换误差 \emph{在
绝对值意义下}相加，放弃了
交换之间的一切抵消。傅里叶证明比较的是特征函数，
其中误差连同其振荡相位一并出现；
埃森光滑化不等式将 $\abs{\varphi_{T_n}
- \varphi_N}$ 以 $\frac{\dd\xi}{\abs\xi}$ 为权积分后，
化为分布函数的界，仅付出对数量级的
代价，并给出贝里--埃森的 $C\beta n^{-1/2}$，
由三阶矩得出。替换以最优性换取
稳健性 --- 正如第 VII 部分所示，还换来可移植性。

\textbf{14.} $\Sigma$ 为对称半正定矩阵；
记 $\Sigma = PDP^{\mathsf T}$（$P$ 正交，$D \geq 0$
为对角阵，见\cref{exo:b3:submanifolds:8}），则对称矩阵 $C =
P\sqrt DP^{\mathsf T}$ 满足 $C^2 = \Sigma$。对任意 $t
\in \R^2$，$\langle t, CZ\rangle = \langle Ct, Z\rangle =
(Ct)_1Z^1 + (Ct)_2Z^2$ 是
独立高斯变量的线性组合，因而服从高斯分布
（见\cref{exo:b3:clt:3}）：$N = CZ$ 为高斯向量；其
均值为 $0$，协方差 $\E[NN^{\mathsf T}] =
C\,\E[ZZ^{\mathsf T}]\,C^{\mathsf T} = CC^{\mathsf T} =
\Sigma$。矩：$\norm N^3 \leq (\abs{N_1} +
\abs{N_2})^3 \leq 4(\abs{N_1}^3 + \abs{N_2}^3)$（由
$x^3$ 在 $\R_+$ 上的凸性），且每个分量都是实
高斯变量，各阶矩均存在
（见\cref{exo:b3:product:10}）：$\gamma' < \infty$。最后
每个 $\langle t, G_n\rangle = \frac1{\sqrt
n}\sum_i\langle t, N_i\rangle$ 为
独立同分布的 $\mathcal N(0, t^{\mathsf T}\Sigma t)$ 变量之标准化和，因而
恰好服从 $\mathcal N(0, t^{\mathsf T}\Sigma t)$：$G_n$ 为
均值为 $0$、协方差为 $\Sigma$ 的高斯向量，且
其分布为 $\mathcal N(0, \Sigma)$
（见\cref{def:b3:clt:gaussianvector}：分布由
这些数据确定）。

\textbf{15.} 设 $\phi(t) = f(w + th)$，$t \in \intcc01$：
$\phi$ 为 $\mathcal C^3$ 类，且
\[
\phi'''(t) = \sum_{j,k,l\in\{1,2\}}\partial_{jkl}f(w +
th)\,h_jh_kh_l, \qquad \abs{\phi'''(t)} \leq
M_3\Bigl(\sum_j\abs{h_j}\Bigr)^3 = M_3(\abs{h_1} +
\abs{h_2})^3 .
\]
取 $3$ 阶泰勒–拉格朗日公式，对 $\phi$ 在 $0$ 与
$1$ 之间，给出第一个不等式；柯西–施瓦茨不等式给出
$\abs{h_1} + \abs{h_2} \leq \sqrt2\norm h$，由此得
常数 $\frac{2\sqrt2M_3}6 = \frac{\sqrt2M_3}3$。

\textbf{16.} 仿第 1 题定义 $H_i$ 与 $W_i$，现
在 $\R^2$ 中，交换论证无需改动。在
第 $i$ 次交换中，一阶项给出
$\sum_j\E[\partial_jf(W_i)]\,(\E X_{i,j} - \E N_{i,j})/
\sqrt n = 0$，二阶项给出
$\frac1{2n}\sum_{j,k}\E[\partial_{jk}f(W_i)]\,(\Sigma_{jk}
- \Sigma_{jk}) = 0$：均值与协方差相符。
由第 15 题给出两个余项的界：
\[
\bigl|\E f(H_i) - \E f(H_{i-1})\bigr| \leq
\frac{\sqrt2M_3}3\cdot\frac{\E\norm{X_i}^3 +
\E\norm{N_i}^3}{n^{3/2}}
= \frac{\sqrt2M_3(\beta' + \gamma')}{3\,n^{3/2}},
\]
并对 $n$ 次交换裂项相消：
\[
\Bigl|\E f\Bigl(\frac{S_n}{\sqrt n}\Bigr) - \E
f(G_n)\Bigr| \leq \frac{\sqrt2\,M_3(\beta' +
\gamma')}{3\sqrt n},
\qquad G_n \sim \mathcal N(0, \Sigma)\ \text{恰好} .
\]
加强论证：$\E\norm{T_n}^2 = \E\norm{X_1}^2 =
\operatorname{tr}\Sigma$（交叉项因
独立性与中心化而消失），故 $\P(\norm{T_n} > A) \leq
\operatorname{tr}\Sigma/A^2$，对 $N$ 同理，
即紧性。对有界连续函数 $f$ 与 $\varepsilon
> 0$，乘以光滑平台函数 $\chi$，它取值 $1$
于半径为 $A$ 的球上，且支集含于半径 $A + 1$ 的球内
（在 $\R^2$ 中对示性函数作光滑化，
见\cref{thm:b3:lp:regularization}）；$g = f\chi$ 一致
连续且具紧支集，故其二维
光滑化 $g_\eta$ 为 $\mathcal C^\infty$ 类，各阶导数
均有界，且 $\norm{g - g_\eta}_\infty
\leq \varepsilon$，只要 $\eta$ 充分小。第 5 题的三 $\varepsilon$
链可原样搬用：$\E f(T_n)
\to \E f(N)$ 对每个有界连续函数 $f \colon \R^2
\to \R$。这即 $d = 2$ 时的\cref{thm:b3:clt:multiclt}，
现已证明——交换法绕开了二维
莱维定理，本章曾将其留作承认。

\textbf{17.} 对有界连续函数 $g \colon \R \to \R$，
映射 $x \mapsto g(\langle t, x\rangle)$ 有界
连续（定义于 $\R^2$ 上），故由第 16 题得 $\E g(\langle
t, T_n\rangle) \to \E g(\langle t, N\rangle)$：每个
投影依分布收敛，且 $\langle t,
N\rangle \sim \mathcal N(0, t^{\mathsf T}\Sigma t)$。（这
是克拉美–沃尔德定理的易证方向：联合收敛
蕴含所有线性像的收敛。）应用：
$V_i = (\xi_i, \xi_i^2 - 1)$ 为独立同分布的中心化向量
（$\E\xi_1^2 = 1$），协方差元素 $\V(\xi_1) =
1$，$\operatorname{Cov}(\xi_1, \xi_1^2 - 1) = \E\xi_1^3$
及 $\V(\xi_1^2 - 1) = \E\xi_1^4 - 1$；三阶矩
$\E\norm{V_1}^3 \leq 4\bigl(\E\abs{\xi_1}^3 +
\E\abs{\xi_1^2 - 1}^3\bigr)$ 有限，当 $\xi_1 \in
L^6$。由第 16 题可得上式所示的联合高斯
极限，并由\cref{thm:b3:clt:gaussianvector}(2)，两个
极限分量独立，当且仅当
协方差 $\E\xi_1^3$ 为零 --- 对于对称分布，
经验均值与经验方差
渐近解耦。

\textbf{18.} 写 $g(x) - g(\theta) = (g'(\theta) +
\eta(x))(x - \theta)$，其中 $\eta(x) = \frac{g(x) -
g(\theta)}{x - \theta} - g'(\theta)$，当 $x \neq \theta$
时，且 $\eta(\theta) = 0$：$\theta$ 处的可导性
恰好等价于 $\eta(x) \to 0$（当 $x \to \theta$ 时）。
\emph{第一步：} $\hat\theta_n \to \theta$ 依概率收敛：
对 $\varepsilon > 0$ 及任意 $A > 0$，从某项起
$\varepsilon\sqrt n \geq A$，故 $\P(\abs{\hat\theta_n -
\theta} > \varepsilon) \leq \P(\abs{\sqrt
n(\hat\theta_n - \theta)} > A) \to \P(\sigma\abs N > A)$
（分布函数在连续点
$\pm A$），且右端趋于 $0$，当 $A \to
\infty$ 时。\emph{第二步：} $\eta(\hat\theta_n) \to 0$ 依
概率：给定 $\varepsilon' > 0$，取 $\delta$ 使得
$\abs\eta \leq \varepsilon'$ 在 $\abs{x - \theta} \leq
\delta$ 上成立；则 $\P(\abs{\eta(\hat\theta_n)} > \varepsilon')
\leq \P(\abs{\hat\theta_n - \theta} > \delta) \to 0$。
\emph{第三步：}
\[
\sqrt n\bigl(g(\hat\theta_n) - g(\theta)\bigr) =
g'(\theta)\,\sqrt n(\hat\theta_n - \theta) +
\eta(\hat\theta_n)\cdot\sqrt n(\hat\theta_n - \theta) .
\]
由斯卢茨基乘积法则（见\cref{exo:b3:clt:8}），取
序列 $\eta(\hat\theta_n) \to 0$ 依概率收敛，且
依分布收敛的 $\sqrt n(\hat\theta_n - \theta)$），则第二项
依分布收敛于 $0\cdot\mathcal N(0,
\sigma^2) = 0$，故依概率收敛于 $0$
（见\cref{exo:b3:clt:4}(b)），第一项依分布收敛于
$g'(\theta)\mathcal N(0, \sigma^2)$（再用斯卢茨基，或
特征函数的仿射规则）；斯卢茨基和
法则将两者合并，极限为 $\mathcal N(0,
g'(\theta)^2\sigma^2)$。

\textbf{19.} (a) 由中心极限定理，$\sqrt n(\bar X_n - \mu)
\Rightarrow \mathcal N(0, \sigma^2)$；由德尔塔方法
取 $g(x) = x^2$，$g'(\mu) = 2\mu$，得 $\sqrt n(\bar
X_n^2 - \mu^2) \Rightarrow \mathcal N(0, 4\mu^2\sigma^2)$
--- 退化（极限为 $0$），当 $\mu = 0$ 时。此时
波动落在高一个量级的尺度上：$n\bar X_n^2 =
(\sqrt n\,\bar X_n)^2$，且对 $t > 0$，
\[
\P\bigl(n\bar X_n^2 \leq t\bigr) = \P\bigl(-\sqrt t \leq
\sqrt n\,\bar X_n \leq \sqrt t\bigr) \longrightarrow
\Phi\Bigl(\frac{\sqrt t}\sigma\Bigr) -
\Phi\Bigl(-\frac{\sqrt t}\sigma\Bigr) = \P(\sigma^2N^2
\leq t) :
\]
$n\bar X_n^2 \Rightarrow \sigma^2N^2$，即高斯变量的平方
（“卡方”型分布） --- 当一阶
导数消失时，泰勒展开的二阶项决定了
非高斯极限。(b) 此处 $\sqrt n(\hat p_n - p)
\Rightarrow \mathcal N(0, p(1 - p))$，且 $g(p) =
\arcsin\sqrt p$ 的导数为 $g'(p) = \frac1{2\sqrt{p(1 - p)}}$，故
$g'(p)^2\,p(1 - p) = \frac14$：极限为 $\mathcal
N(0, \frac14)$，对所有 $p \in \intoo01$ 成立。在
$\arcsin$ 尺度上，渐近 $95\%$ 误差界为 $\pm
\frac{0.98}{\sqrt n}$，事先已知 --- 而在
\cref{ex:b3:clt:confidence}中，宽度涉及未知的
$\sigma = \sqrt{p(1-p)}$，需按最坏情形取 $\frac12$
或加以估计：该变换\emph{稳定}
了方差。

\textbf{20.} 令 $A^* = \{k : \mu(\{k\}) > \nu(\{k\})\}$
且 $\Delta_k = \mu(\{k\}) - \nu(\{k\})$，于是
$\sum_k\Delta_k = 0$。对任意 $A \subseteq \N$：$\mu(A) -
\nu(A) = \sum_{k\in A}\Delta_k \leq \sum_{k\in
A^*}\Delta_k$，等号在 $A = A^*$ 时成立；又因
$(\Delta_k)$ 的正部与负部
总质量相等，$\sum_{A^*}\Delta_k =
\frac12\sum_k\abs{\Delta_k}$。交换 $\mu, \nu$ 即可消去符号，
即得 $d_{\mathrm{TV}}(\mu, \nu) =
\frac12\sum_k\abs{\Delta_k}$。耦合：对任意 $A$，
\[
\mu(A) - \nu(A) = \E\bigl[\mathbf 1_A(X) - \mathbf
1_A(Y)\bigr] = \E\bigl[(\mathbf 1_A(X) - \mathbf
1_A(Y))\,\mathbf 1_{X\neq Y}\bigr] \leq \P(X \neq Y),
\]
再对 $A$ 取上确界。

\textbf{21.} 两分布的质量分别为：$k = 0$ 时 $1 - p$ 对
$\eu^{-p}$，且 $\eu^{-p} > 1 - p$；$k = 1$ 时 $p$ 对
$p\,\eu^{-p} < p$；$k \geq 2$ 时 $0$ 对泊松
余项 $1 - \eu^{-p} - p\eu^{-p} \geq 0$。于是
\[
d_{\mathrm{TV}} = \tfrac12\bigl[(\eu^{-p} - 1 + p) + (p -
p\eu^{-p}) + (1 - \eu^{-p} - p\eu^{-p})\bigr] =
\tfrac12\bigl(2p - 2p\eu^{-p}\bigr) = p(1 - \eu^{-p}),
\]
且 $1 - \eu^{-p} \leq p$，给出界 $p^2$。

\textbf{22.} 记 $H_{i-1} = W_i + X_i$ 及 $H_i = W_i +
Y_i$，其中 $W_i = \sum_{j<i}Y_j + \sum_{j>i}X_j$，
与配对 $(X_i, Y_i)$ 独立（按组划分）。对 $A
\subseteq \N$，对可数多个取值作条件化，
由独立性，
\[
\P(H_{i-1} \in A) = \sum_{k\geq0}\P(X_i = k)\,\P(W_i + k
\in A),
\]
对 $H_i$ 同理，代以 $Y_i$。相减，记 $c_k
= \P(W_i + k \in A) \in \intcc01$，且 $\Delta_k = \P(X_i =
k) - \P(Y_i = k)$，其和为零：
\[
\abs{\P(H_{i-1} \in A) - \P(H_i \in A)} =
\Bigl|\sum_k\Delta_k\bigl(c_k - \tfrac12\bigr)\Bigr| \leq
\tfrac12\sum_k\abs{\Delta_k} =
d_{\mathrm{TV}}\bigl(\mathcal B(1, p_i), \mathcal
P(p_i)\bigr) .
\]
逐项相消，从 $H_0 = S$ 到 $H_n = \sum_iY_i \sim
\mathcal P(\lambda)$（见\cref{exo:b3:clt:1}，反复应用），并
利用第 21 题：
\[
\abs{\P(S \in A) - \P(\mathcal P(\lambda) \in A)} \leq
\sum_{i=1}^np_i\bigl(1 - \eu^{-p_i}\bigr) \leq
\sum_{i=1}^np_i^2
\qquad\text{ 对每个 } A :
\]
勒卡姆不等式。（第 20 题的耦合界给出另一
种证法：每对在同一均匀随机变量上耦合，
使 $\P(X_i \neq Y_i) \leq p_i^2$，并估计
$\P(S \neq \sum Y_i)$；替换本身无需任何构造
。）

\textbf{23.} (a) 取 $p_i = \frac\lambda n$：
$d_{\mathrm{TV}}(\text{分布 }S, \mathcal P(\lambda))
\leq \frac{\lambda^2}n$。这在
\cref{exo:b3:clt:5}的基础上加强了三点：其一，在每个有限的 $n$ 上给出显式误差，
其二，对所有事件 $A$ 一致成立，
（而非逐个区间），其三，不必要求
$p_i$ 相等 --- 只需 $\sum_ip_i^2$ 小，例如 $\sum p_i^2 \leq
\lambda\max_ip_i$：\emph{稀有事件众多，
无一占主导}。(b) 此处 $n = 500$，$p_i = \frac1{500}$，
$\lambda = 1$：泊松模型的误差至多为
$500\cdot\frac1{500^2} = 0.002$（对每个事件）；特别地，
取 $A = \{0\}$，
\[
\P(\text{无信件错投}) = \Bigl(1 -
\frac1{500}\Bigr)^{500},
\qquad
\Bigl|\P(\text{无信件错投}) - \eu^{-1}\Bigr| \leq
0.002,
\]
于是所求值为 $\eu^{-1} \approx 0.368$，
可保证误差不超过 $0.002$（真实偏差约为
$4\cdot10^{-4}$）。(c) 本题
归结为一种方法，有两种情形。当 $n$
个同量级的贡献各自的方差为
$\frac1n$，此时与高斯分布匹配\emph{前两}阶矩，
使每次替换的误差为 $o(\frac1n)$，于是和趋于
高斯 --- 以泰勒展开作为局部比较工具。
当 $n$ 个贡献都是概率为
$p_i$ 的示性变量时，将\emph{均值}与泊松点质量匹配，
使每次替换的代价为 $p_i^2$：稀有事件的计数趋于
泊松分布 --- 以全变差作为精确的局部
比较工具。同样的混合、同样的逐项相消，
局部估计各不相同：替换是一种策略，而非定理，
高斯与泊松极限是它最古老的两项
收获。

\textbf{24.} 由中心极限定理，$\sqrt n(\bar X_n - \mu)
\Rightarrow \mathcal N(0, \sigma^2)$，且 $g(x) = \ln x$ 在
$\mu > 0$ 处可导，且 $g'(\mu) = \frac1\mu$，由
delta 方法（第 VI 部分）可得 $\sqrt n(\ln\bar X_n - \ln\mu)
\Rightarrow \mathcal N(0, \sigma^2/\mu^2)$。经指数化还原
区间 $\abs{\ln\bar X_n - \ln\mu} \leq
\frac{1.96\,\sigma}{\mu\sqrt n}$，取指数即得：
\[
\mu \in \bar X_n\cdot
\eu^{\pm1.96\,\sigma/(\mu\sqrt n)}
\qquad\text{渐近概率为 } 95\%
\]
（实际中用 $\sigma/\mu$ 的经验版本替换之，
并借助斯卢茨基定理，如\cref{exo:b3:clt:8}所示）。
乘性区间是自然的选择，适用于数据
为正，且误差与其大小成正比 ---
收入、浓度、半衰期：这些量天然
处于对数尺度，对称的加性区间可能
甚至跨过零点。

\textbf{25.} $\E[(X - p)^3] = (1-p)^3p + (-p)^3(1 - p) =
p(1-p)\bigl[(1-p)^2 - p^2\bigr] = p(1-p)(1 - 2p)$。在
交换分析（第 IV 部分）中，两阶矩匹配之后的
主误差项携带的是\emph{带符号的}
三阶矩：当 $p < \frac12$ 时为正（分布向
右偏：偶有远高于均值的大幅偏离），而
正态近似系统性地错置概率质量 ---
低估较短的左尾，高估
右尾 --- 误差阶为 $n^{-1/2}$；当 $p =
\frac12$ 时三阶矩为零，伯努利分布与
高斯分布在三阶上吻合，收敛速度随之提高（第
IV 部分的矩匹配题）。数值上：$\P(S = 0) =
0.9^{20} = 0.1216$，而高斯分布
$\mathcal N(2, 1.8)$ 给出 $\Phi\bigl(\frac{0.5 -
2}{\sqrt{1.8}}\bigr) = \Phi(-1.118) \approx 0.132$：
正态曲线不知 $0$ 处的边界，也不知
向右偏斜，在底部堆积了过多的概率质量 ---
这正是预言的误差符号，在 $n = 20$ 时即可见。
\end{solution}
